\documentclass[10pt,a4paper]{amsart}
\usepackage{amsmath,amssymb,amsthm,graphicx}
\usepackage{verbatim}
\newenvironment{reference}{\comment}{\endcomment}
%\newenvironment{reference}{}{}
\newenvironment{slogan}{\comment}{\endcomment}
\newenvironment{history}{\comment}{\endcomment}
% For commutative diagrams we use Xy-pic
\usepackage[all]{xy}

% We use 2cell for 2-commutative diagrams.
\xyoption{2cell}
% Theorem environments.
%
\theoremstyle{plain}
\newtheorem{theorem}[subsection]{Theorem}
\newtheorem{proposition}[subsection]{Proposition}
\newtheorem{lemma}[subsection]{Lemma}

\theoremstyle{definition}
\newtheorem{definition}[subsection]{Definition}
\newtheorem{example}[subsection]{Example}
\newtheorem{exercise}[subsection]{Exercise}
\newtheorem{situation}[subsection]{Situation}

\theoremstyle{remark}
\newtheorem{remark}[subsection]{Remark}
\newtheorem{remarks}[subsection]{Remarks}

\numberwithin{equation}{subsection}

% Macros
%
\def\lim{\mathop{\mathrm{lim}}\nolimits}
\def\colim{\mathop{\mathrm{colim}}\nolimits}
\def\Spec{\mathop{\mathrm{Spec}}}
\def\Hom{\mathop{\mathrm{Hom}}\nolimits}
\def\Ext{\mathop{\mathrm{Ext}}\nolimits}
\def\SheafHom{\mathop{\mathcal{H}\!\mathit{om}}\nolimits}
\def\SheafExt{\mathop{\mathcal{E}\!\mathit{xt}}\nolimits}
\def\Sch{\mathit{Sch}}
\def\Mor{\mathop{\mathrm{Mor}}\nolimits}
\def\Ob{\mathop{\mathrm{Ob}}\nolimits}
\def\Sh{\mathop{\mathit{Sh}}\nolimits}
\def\NL{\mathop{N\!L}\nolimits}
\def\CH{\mathop{\mathrm{CH}}\nolimits}
\def\proetale{{pro\text{-}\acute{e}tale}}
\def\etale{{\acute{e}tale}}
\def\QCoh{\mathit{QCoh}}
\def\Ker{\mathop{\mathrm{Ker}}}
\def\Im{\mathop{\mathrm{Im}}}
\def\Coker{\mathop{\mathrm{Coker}}}
\def\Coim{\mathop{\mathrm{Coim}}}

% Boxtimes
%
\DeclareMathSymbol{\boxtimes}{\mathbin}{AMSa}{"02}

%
% Macros for moduli stacks/spaces
%
\def\QCohstack{\mathcal{QC}\!\mathit{oh}}
\def\Cohstack{\mathcal{C}\!\mathit{oh}}
\def\Spacesstack{\mathcal{S}\!\mathit{paces}}
\def\Quotfunctor{\mathrm{Quot}}
\def\Hilbfunctor{\mathrm{Hilb}}
\def\Curvesstack{\mathcal{C}\!\mathit{urves}}
\def\Polarizedstack{\mathcal{P}\!\mathit{olarized}}
\def\Complexesstack{\mathcal{C}\!\mathit{omplexes}}
% \Pic is the operator that assigns to X its picard group, usage \Pic(X)
% \Picardstack_{X/B} denotes the Picard stack of X over B
% \Picardfunctor_{X/B} denotes the Picard functor of X over B
\def\Pic{\mathop{\mathrm{Pic}}\nolimits}
\def\Picardstack{\mathcal{P}\!\mathit{ic}}
\def\Picardfunctor{\mathrm{Pic}}
\def\Deformationcategory{\mathcal{D}\!\mathit{ef}}
\title{The unbounded derived category of modules supported in a specialization-stable subset is the full subcategory of complexes supported there, in finite dimension}
\author{The Stacks Project Authors}
\date{stacks-project:local-cohomology.tex:local-snapshot-20261012; GNU FDL 1.2 or later}
\begin{document}
\maketitle
\noindent\textit{Editorial note.} All delivered mathematical text is the author's own, copied from the pinned Stacks Project chapter source sources/stacks-local-cohomology/pinned-original-source.tex; no statement or proof step was rewritten, reordered or removed.\par\medskip
\noindent\textit{Editorial note.} Preamble provenance: the Stacks Project keeps one single shared preamble for every chapter; this card takes it directly from the staged sources/stacks-local-cohomology/preamble.tex as three verbatim spans (lines 8-12, 14-18 and 154-217), so no derived file was written and nothing was added or redefined. The wrapper class and the base packages are supplied by the assembler.\par\medskip
\noindent\textit{Editorial note.} Bibliography provenance: neither the counted statement nor its proof contains a citation (this lane rescans the whole span before writing the card), so the card carries no bibliography and no entry was reconstructed.\par\medskip
\noindent\textit{Editorial note.} Selection provenance: this card was opened under the file-based control contract (CONTROL.md), with the mechanical pre-card gate run on its statement and proof interval before the card was written.\par\medskip
\noindent\textit{Allowed standard prerequisites (editorial).}\begin{itemize}
\item The vanishing of local cohomology with support in a closed subset above the dimension of the ring. Author lemma lemma-cd-dimension, whose statement is reproduced verbatim as a context span of this card (source lines 758-764). Gives the bound on the cohomological dimension of the support functor.
\item The colimit description of the sections with support in a specialization-stable subset, and the bounded below case of the equivalence. Author lemmas lemma-adjoint-ext and lemma-equal-plus, whose statements are reproduced verbatim as context spans of this card (source lines 996-1004 and 1019-1024). The first transfers the vanishing of the local cohomology to the general support, the second supplies the equivalence on bounded below complexes.
\end{itemize}\medskip
\begin{lemma}
\label{lemma-equal-full}
Let $A$ be a Noetherian ring. Let $T \subset \Spec(A)$ be a subset stable
under specialization. If $\dim(A) < \infty$, then functor
$D(\text{Mod}_{A, T}) \to D_T(A)$ is an equivalence.
\end{lemma}

\begin{lemma}
\label{lemma-cd-dimension}
Let $I \subset A$ be a finitely generated ideal of a ring $A$.
If $M$ is a finite $A$-module, then
$H^i_{V(I)}(M) = 0$ for $i > \dim(\text{Supp}(M))$.
In particular, we have $\text{cd}(A, I) \leq \dim(A)$.
\end{lemma}

\begin{lemma}
\label{lemma-adjoint-ext}
Let $A$ be a Noetherian ring. Let $T \subset \Spec(A)$
be a subset stable under specialization.
For any object $K$ of $D(A)$ we have
$$
H^i(RH^0_T(K)) = \colim_{Z \subset T\text{ closed}} H^i_Z(K)
$$
\end{lemma}

\begin{lemma}
\label{lemma-equal-plus}
Let $A$ be a Noetherian ring. Let $T \subset \Spec(A)$ be a subset stable
under specialization. The functor $D^+(\text{Mod}_{A, T}) \to D^+_T(A)$
is an equivalence.
\end{lemma}

\begin{proof}
Say $\dim(A) = d$. Then we see that $H^i_Z(M) = 0$ for $i > d$
for every closed subset $Z$ of $\Spec(A)$, see
Lemma \ref{lemma-cd-dimension}.
By Lemma \ref{lemma-adjoint-ext} we find that $H^0_T$ has bounded
cohomological dimension.

\medskip\noindent
Let $K \in D_T(A)$. We claim that $RH^0_T(K) \to K$ is an
isomorphism. We know this is true when $K$ is bounded below, see
Lemma \ref{lemma-equal-plus}. However, since $H^0_T$ has bounded
cohomological dimension, we see that the $i$th cohomology of
$RH_T^0(K)$ only depends on $\tau_{\geq -d + i}K$ and we conclude.
Thus $D(\text{Mod}_{A, T}) \to D_T(A)$ is an equivalence with
quasi-inverse $RH^0_T$.
\end{proof}

\end{document}
